\chapter{Лабораторная работа №2. Создание взаимосвязанной IP-сети}
\label{ch:chapter 2}

\section{\textit{Лабораторная работа №2.1. Адресация и маршрутизация IPv4}} \label{sec:lab2.1}

\subsection{Цели}

Лабораторная работа помогает получить практические навыки по изучению
следующих тем:

\begin{itemize}

\item Процедура настройки IPv4-адреса на интерфейсе

\item Функции и значение loopback-интерфейсов

\item Принципы генерирования прямых маршрутов

\item Процедура настройки статических маршрутов и условия, при которых используются статические маршруты

\item Процедура проверки возможности установления соединения сетевого уровня с помощью инструмента ping

\item Процедура настройки статических маршрутов и сценарии их применения

\end{itemize}

\subsection{Топология сети}

Ниже представлена топология сети для данной работы.

\includegraphics[width=0.5\textwidth]{lb2_1.png} 
\captionof{figure}{Топология сети для лабораторной работы №2}

\subsection{План работы}

\begin{itemize}

\item 1. Настройка IP-адресов для интерфейсов на маршрутизаторах.


\item 2. Настройка статических маршрутов для установления связи между маршрутизаторами.

\end{itemize}

\subsection{Процедура конфигурирования}

\textbf{Шаг 1.} Настроим основные параметры устройства. 

Зададим им имена через команду \textbf{sysname}

\includegraphics[width=0.5\textwidth]{lb2_2.png} 

\includegraphics[width=0.5\textwidth]{lb2_3.png} 

\includegraphics[width=0.5\textwidth]{lb2_4.png} 

\textbf{Шаг 2.} Выведем на экран статус интерфейса на маршрутизаторах с помощью команды \textbf{display ip interface brief}.

\includegraphics[width=0.5\textwidth]{lb2_5.png} 

\includegraphics[width=0.5\textwidth]{lb2_6.png} 

\includegraphics[width=0.5\textwidth]{lb2_7.png} 

Теперь выведем на экран таблицу маршрутизации на маршрутизаторах с помощью команды \textbf{display ip routing-table}:

\includegraphics[width=0.5\textwidth]{lb2_8.png} 

\includegraphics[width=0.5\textwidth]{lb2_9.png} 

\includegraphics[width=0.5\textwidth]{lb2_10.png} 

\textbf{Шаг 3.} Настройка IP-адресов.

Настроим IP-адреса для физических интерфейсов в соответствии с таблицей:

\includegraphics[width=0.5\textwidth]{lb2_11.png} 

\includegraphics[width=0.5\textwidth]{lb2_12.png}\hfill\includegraphics[width=0.5\textwidth]{lb2_13.png} 

\includegraphics[width=0.5\textwidth]{lb2_14.png}\hfill\includegraphics[width=0.5\textwidth]{lb2_15.png} 

\includegraphics[width=0.5\textwidth]{lb2_16.png}\hfill\includegraphics[width=0.5\textwidth]{lb2_17.png}

Проверим связь роутеров через команду \textbf{ping}:

Проверим связность R1 с R2 и R3:

\includegraphics[width=0.5\textwidth]{lb2_21.png}

Как видно, связность присутствует.
 
Теперь проверим настройку интерфейсов при помощи команды \textbf{display ip interface brief}

\includegraphics[width=0.5\textwidth]{lb2_18.png}

\includegraphics[width=0.5\textwidth]{lb2_19.png}

\includegraphics[width=0.5\textwidth]{lb2_20.png}

\textbf{Шаг 4.} Создание loopback-интерфейса.

Настраиваться он будет в соответствии со следующей таблицей:

\includegraphics[width=0.5\textwidth]{lb2_22.png}

Loopback-интерфейсы — это настроенные вручную логические интерфейсы, которые
физически не существуют. Логические интерфейсы могут использоваться для обмена
данными. Loopback-интерфейс всегда находится в рабочем состоянии (статус Up) на
физическом и канальном уровнях, если только он не был отключен вручную. Обычно
loopback-интерфейс имеет 32-битную маску. Loopback-интерфейсы используются в
следующих случаях:

\begin{itemize}
    \item 1.В качестве адреса для идентификации и управления маршрутизатором .
    \item 2.В качестве идентификатора маршрутизатора в OSPF.
    \item 3.Для повышения надежности сети.
\end{itemize}

В данной работе loopback-интерфейсы используются для имитации клиентов

Чтобы назначить адреса на loopback-интерфейсы, воспльзуемся командой \textbf{interface LoopBacko}:

\includegraphics[width=0.5\textwidth]{lb2_24.png}

\includegraphics[width=0.5\textwidth]{lb2_25.png}

\includegraphics[width=0.5\textwidth]{lb2_26.png}

Выведем на экран таблицу маршрутизации для R1:

\includegraphics[width=0.5\textwidth]{lb2_23.png}

Проверим наличие связности между loopback-интерфейсами:

\includegraphics[width=0.5\textwidth]{lb2_27.png}

Связности, как видно, нет, т.к. мы не настроили маршрут между этими интерфейсами.

\textbf{Шаг 4.} Настройка статических маршрутов.

На R1 настроим маршруты к LoopBack 0 R2 и R3:

\includegraphics[width=0.5\textwidth]{lb2_28.png}

Теперь выведем таблицу маршрутизации:

\includegraphics[width=0.5\textwidth]{lb2_29.png}

Теперь снова проверим связность:


\includegraphics[width=0.5\textwidth]{lb2_30.png}

Снова не работает, т.к. маршрут не настроен для R2:

\includegraphics[width=0.5\textwidth]{lb2_31.png}

Снова проверим связность:

\includegraphics[width=0.5\textwidth]{lb2_32.png}

Теперь связность присутствует

По аналогии настроим маршруты на других роутерах и проверим связность:

\includegraphics[width=0.5\textwidth]{lb2_33.png}

\includegraphics[width=0.5\textwidth]{lb2_34.png}

\includegraphics[width=0.5\textwidth]{lb2_35.png}

\includegraphics[width=0.5\textwidth]{lb2_36.png}

\textbf{Шаг 6.} Настроим маршрут от R1 к R2 через R3 в качестве резервного маршрута от LoopBack0
R1 к LoopBack0 R2.

Для этого настроим статические маршруты на R1 и R2:

\includegraphics[width=0.5\textwidth]{lb2_37.png}

\includegraphics[width=0.5\textwidth]{lb2_38.png}

Выведем таблицы маршрутизации для R1 и R2:

\includegraphics[width=0.5\textwidth]{lb2_39.png}

\includegraphics[width=0.5\textwidth]{lb2_40.png}

Отключим интерфейс GigabitEthernet0/0/3 на маршрутизаторах R1 и R2, чтобы
сделать недействительным маршрут с наивысшим приоритетом:

\includegraphics[width=0.5\textwidth]{lb2_41.png}

Теперь выведем таблицу маршрутизации, чтобы проверить, включился ли другой LoopBack интерфейс:

\includegraphics[width=0.5\textwidth]{lb2_42.png}

Проверим установление связей:

\includegraphics[width=0.5\textwidth]{lb2_43.png}

Как видим, она присутствует

Теперь проверим, точно ли через запасной маршрут мы устанавливаем связность с помощью \textbf{tracert}:

\includegraphics[width=0.5\textwidth]{lb2_44.png}

\textbf{Шаг 7.} Настроим маршруты по умолчанию для установления связи между интерфейсом
LoopBack0 маршрутизатора R1 и интерфейсом LoopBack0 маршрутизатора R2.

Включим выключенные интерфейсы и удалим настроенные маршруты:

\includegraphics[width=0.5\textwidth]{lb2_45.png}

Выведем таблицу маршрутизации:

\includegraphics[width=0.5\textwidth]{lb2_46.png}

Настроим маршрут по умолчанию для R1:

\includegraphics[width=0.5\textwidth]{lb2_47.png}

Выведем таблицу маршрутизации:

\includegraphics[width=0.5\textwidth]{lb2_48.png}

Проверим наличие связности между loopback-интерфейсами R1 и R2:

\includegraphics[width=0.5\textwidth]{lb2_49.png}   

Связность присутствует.

\section{Вопросы}

1. В каких ситуациях настроенный статический маршрут будет добавлен в таблицу
IP-маршрутизации? Можно ли добавить маршрут в таблицу IP-маршрутизации,
если настроенный следующий переход недоступен? 

Статический маршрут переходит в активное состояние и добавляется в таблицу маршрутизации (\textit{Routing Table}) только при выполнении следующих условий:
\begin{itemize}
    \item \textbf{Статус интерфейса:} Выходной интерфейс должен находиться в состоянии \texttt{Up/Up}.
    \item \textbf{Достижимость Next-hop:} Если указан IP-адрес следующего перехода, маршрутизатор должен иметь в таблице активный маршрут к этой подсети (прямое соединение или рекурсивный поиск).
    \item \textbf{Административное расстояние:} Статический маршрут должен иметь меньшее значение Administrative Distance (AD), чем другие маршруты к этой же сети от протоколов динамической маршрутизации.
\end{itemize}

По умолчанию — \textbf{нет}. Если адрес следующего перехода недоступен, маршрут помечается как \textit{Inactive} и не попадает в таблицу продвижения пакетов.

2. Каким будет IP-адрес источника пакетов ICMP на шаге 3, если при проверке связи
между loopback-интерфейсами не указать аргумент -a? Почему?

При проверке связи между Loopback-интерфейсами без аргумента \texttt{-a} (или \texttt{source}), IP-адресом источника будет \textbf{основной IP-адрес исходящего физического интерфейса}.

\textbf{Обоснование:}
\begin{enumerate}
    \item Маршрутизатор выполняет поиск пути к цели в таблице маршрутизации.
    \item После определения выходного интерфейса стек TCP/IP автоматически подставляет его IP-адрес в поле \textit{Source IP} пакета, если иное не задано администратором вручную.
\end{enumerate}

\section{\textit{Лабораторная работа №2.2. Адресация и маршрутизация IPv4}}

\subsection{Цели}

Лабораторная работа помогает получить практические навыки по изучению
следующих тем:

\begin{itemize}

\item Основные команды OSPF

\item Процедура проверки рабочего статуса OSPF

\item Процедура настройки выбора маршрутов OSPF на основании их стоимости

\item Анонсирование маршрутов по умолчанию в OSPF

\item Процедура настройки аутентификации OSPF

\end{itemize}

\subsection{Топология сети}

Ниже представлена топология сети для данной работы.

\includegraphics[width=0.5\textwidth]{lb2_1.png} 
\captionof{figure}{Топология сети для лабораторной работы №2}

\subsection{План работы}

\begin{itemize}

\item 1. Создание процессов OSPF на устройствах и включение OSPF на интерфейсах.


\item 2. Настройка аутентификации OSPF.

\item 3. Настройка OSPF для анонсирования маршрутов по умолчанию.

\item 4. Управление выбором маршрутов OSPF на основании их стоимости.

\end{itemize}

\subsection{Процедура конфигурирования}

\textbf{Шаг 1.} Настроим основные параметры.

Для этого, для начала, нам нужно повторить шаги с 1го по 4й в "Лабораторной работе №\ref{sec:lab2.1}"

Выведем таблицу маршрутизации R1:

\includegraphics[width=0.5\textwidth]{lb2_50.png}   

\textbf{Шаг 2.} Настроим основные параметры OSPF.

Для начала нужно создать процесс OSPF:

\includegraphics[width=0.5\textwidth]{lb2_51.png}   

Настройка параметров OSPF станет возможной только после создания процесса
OSPF. OSPF поддерживает несколько независимых процессов на одном устройстве.
Обмен маршрутами между различными процессами OSPF осуществляется аналогично
обмену маршрутами между разными протоколами маршрутизации. При создании
процесса OSPF можно указать идентификатор процесса. Если идентификатор
процесса не указан, то по умолчанию используется идентификатор процесса 1

Создадим область OSPF и укажем интерфейсы, которые будут им пользоваться:

\includegraphics[width=0.5\textwidth]{lb2_52.png}   

С помощью команды area можно создать область OSPF и перейти в режим
конфигурирования области OSPF.

\includegraphics[width=0.5\textwidth]{lb2_53.png}   

С помощью команды network network-address wildcard-mask можно указать
интерфейсы, на которых необходимо включить OSPF. OSPF будет работать на
интерфейсе только при соблюдении следующих двух условий:

1.
Длина маски IP-адреса интерфейса должна быть не короче, чем длина маски,
указанная в команде network. Для OSPF должна использоваться обратная маска.
Например, 0.0.0.255 указывает, что длина маски составляет 24 бита.

2.
Адрес интерфейса должен находиться в пределах сетевого диапазона,
указанного в команде network.

В данном примере OSPF можно включить на трех интерфейсах, и все они добавлены в
область 0

\includegraphics[width=0.5\textwidth]{lb2_54.png}   

Если обратная маска в команде network включает только нули (0) и IP-адрес
интерфейса совпадает с IP-адресом, указанным в команде network-address, то на
интерфейс также будет работать OSPF.

\includegraphics[width=0.5\textwidth]{lb2_55.png}   

\textbf{Шаг 3.} Выведем на экран рабочий статус OSPF.

Сначала информацию о соседях OSPF:

\includegraphics[width=0.5\textwidth]{lb2_56.png} 

Команда display ospf peer позволяет вывести на экран информацию о соседях в
каждой области OSPF. Информация включает в себя область, к которой принадлежит
сосед, идентификатор маршрутизатора соседа, статус соседа, DR и BDR.

Выведем маршруты, полученные от OSPF.

\includegraphics[width=0.5\textwidth]{lb2_57.png} 

\textbf{Шаг 4.} Настройка аутентификации OSPF:

Сначала на R1:

\includegraphics[width=0.5\textwidth]{lb2_58.png} 

При просмотре конфигурации пароль отображается в зашифрованном виде,
поскольку в команде указано слово «cipher», обеспечивающее шифрование текста.

Выведем на экран OSPF соседей:

\includegraphics[width=0.5\textwidth]{lb2_59.png} 

Теперь настроим аутентификацию на R2:

\includegraphics[width=0.5\textwidth]{lb2_60.png} 

Выведем соседей OSPF, теперь уже на R2:

\includegraphics[width=0.5\textwidth]{lb2_61.png}

Теперь мы видим соседа в лице R1, т.к. аутентификация настроена на R1 и R2.

Теперь настроим аутентификацию на R3:

\includegraphics[width=0.5\textwidth]{lb2_62.png}

\textbf{Шаг 5.} Пусть R1 - граничный роутер всех сетей. Тогда он анонсирует маршрут OSPF по умолчанию.

Анонсируем его:

\includegraphics[width=0.5\textwidth]{lb2_63.png}

Команда default-route-advertise позволяет анонсировать маршрут по умолчанию в
общую область OSPF. Если аргумент always не указан, маршрут по умолчанию
анонсируется другим маршрутизаторам только тогда, когда в таблице
маршрутизации локального маршрутизатора есть активные маршруты по умолчанию
других протоколов, не OSPF. В данном случае в локальной таблице маршрутизации
маршрут по умолчанию отсутствует. Таким образом, необходимо использовать
аргумент always.

Выведем таблицы маршрутизации R2 и R3:

\includegraphics[width=0.5\textwidth]{lb2_64.png}

\includegraphics[width=0.5\textwidth]{lb2_65.png}

\textbf{Шаг 6.} Настройка аутентификации OSPF.

Изменим значения стоимости интерфейсов на R1, чтобы LoopBack0 на R1 мог
достигать LoopBack0 на R2 через R3.

Согласно таблице маршрутизации R1 стоимость маршрута от маршрутизатора R1 до
LoopBack0 маршрутизатора R2 равна 1, а стоимость маршрута от R1 к R2 через R3
равна 2 Следовательно, необходимо только установить для стоимости маршрута от
маршрутизатора R1 до LoopBack0 маршрутизатора R2 значение больше 2

\includegraphics[width=0.5\textwidth]{lb2_66.png}

Проверим результат с помощью \textbf{tracert}:

\includegraphics[width=0.5\textwidth]{lb2_67.png}

\section{Вопросы.}

Какой маршрут будет использоваться на шаге 6 маршрутизатором R2 для
возврата пакетов ICMP на маршрутизатор R1? Попробуйте объяснить причину.

В OSPF (и маршрутизации вообще) выбор пути происходит на основе локальной таблицы маршрутизации конкретного устройства. Вот почему R2 ответит именно так:
Асимметричная маршрутизация: То, что вы изменили стоимость интерфейса на R1, влияет только на то, как пакеты уходят с R1. Это никак не меняет настройки интерфейсов на самом R2.

Стоимость (Metric):

Для R2 стоимость пути к LoopBack0 R1 через прямое соединение (R2 -> R1) по-прежнему равна 1.
Стоимость пути к R1 через посредника (R2 -> R3 -> R1) равна 2.
Логика роутера: Маршрутизатор всегда выбирает путь с наименьшей стоимостью. Поскольку на интерфейсах R2 вы ничего не меняли, для него «кратчайшим» путем до R1 остается прямая линк-связь.
